\begin{definition}[The implementation of a protocol]
  \label{def:aba-implementation}
  \lean{PLTS.ABA.Implementation.NetworkEvent}\lean{PLTS.ABA.Implementation.actsAt}\lean{PLTS.ABA.Implementation.coinLabelMap}\lean{PLTS.ABA.Implementation.RoundVariablesMap}\lean{PLTS.ABA.Implementation.IsRoundVariables}\lean{PLTS.ABA.Implementation.NetworkState}\lean{PLTS.ABA.Implementation.roundOf}\lean{PLTS.ABA.Implementation.NetworkState.writeGhost}\lean{PLTS.ABA.Implementation.roundOwn}\lean{PLTS.ABA.Implementation.ProgramStep}\lean{PLTS.ABA.Implementation.NetworkStep}\lean{PLTS.ABA.Implementation.IsRoundStep}\lean{PLTS.ABA.Implementation.system}\leanok
  \uses{def:aba-labels, def:aba-parameters, def:aba-round-loop-state, def:relabel, def:synchronised-product, def:process-algebra, def:idle-family}
  A protocol as it runs: $n$ programs, one per process, beside one network and the common coin.
  A program reads its own variables, its own received sets and its own replacement flag, and nothing else about corruption --- not the corrupted set, not the budget, not another process's status (D23) --- while the adversary holds the round-tagged message sets, the DECIDED sets and the corrupted set with its budget, and it alone has transitions on the Byzantine labels.
  The system is parametric in five things a graded-agreement implementation supplies: its round message type $M$, its round's own events $E$, its per-process per-round variables $S$, the payload its graded-agreement call multicasts, if any, and its transitions, given as a relation embedded in one constructor of the program's transitions.
  What does not depend on them is written once --- the round loop, the DECIDED relay and its delivery, the coin's call and return, corruption, the adversary's transitions, the composition pipeline
  \[
    \mathrm{system} \;=\;
      \mathrm{abstract}_{\mathrm{hiddenAPI}}\Bigl(
        \mathrm{relabel}\bigl(\mathrm{abstract}_{\mathrm{NetworkEvent}}(
          (\mathrm{synchronisedProduct}_{j \in [n]}\ \mathrm{program}(j)) \parallel (\mathrm{network} \parallel \mathrm{coinOverExtendedAlphabet}))\bigr)\Bigr),
  \]
  and the lemmas that read a transition off its label.
  Which transitions belong to the implementation is settled by the label: a program's transition is the implementation's business exactly when its label is one of $\mathrm{roundOwn}(j)$ --- the graded-agreement call and return at $j$, $j$'s own round multicast, a delivery addressed to $j$, $j$'s own round events, and $j$'s own call against round variables that already hold an input.
  An implementation must supply four things about its transitions: that they carry such a label, that they fire only at an unreplaced program, that they are Dirac, and that a return transition takes the announced bit free, firing at every other bit with the same successor (D29).
  That is \leandecl{PLTS.ABA.Implementation.IsRoundStep}; the lemmas that read a transition off its label consume the first three, and Theorem~\ref{thm:aba-ghost-removal} the fourth.
  The alphabet the components synchronise on is parametric in $M$ in its round multicast and its round delivery and in $E$ in its round event, a hidden synchronisation on which the network moves and sends nothing, and fixed in every other constructor, and is hidden before it leaves the composition, so two instantiations at the same protocol need not agree on it.
  The adversary holds one thing more, the ghost, and an implementation supplies it too (D30).
  For each round $r$ its state carries a ghost $ghost(r)$ of a type $G$ the implementation fixes.
  The ghost belongs to the network: no program's transition reads it and no program's variables hold it.
  Two parameters carry it.
  The first, $\mathrm{ghostStep}$, writes it: on every transition the ghost of the round the label names --- $\mathrm{roundOf}$ of that label --- is replaced by $\mathrm{ghostStep}$ of the label, the adversary's state and the ghost standing there, the ghosts of the other rounds are left where they stand, and a label naming no round leaves the whole ghost alone.
  That is $\mathrm{writeGhost}$, and it sits on the successor of every transition of the adversary.
  The second, $\mathrm{ghostOutput}$, reads it out: it is a relation between the adversary's state, the round, the process being answered, the graded outcome and a bit, and it guards the two graded-agreement return transitions, the correct $\mathrm{retG}$ and the Byzantine $\mathrm{byzantineRetG}$ at a replaced program, each of which fires only with the bound bit its label carries standing in the relation at the state before the transition.
  An implementation that computes the announced bit writes the ghost at the first return of a round, from the sent sets and the corrupted set, reads the same ghost back at every later return of that round, and instantiates the relation as the equation between the announced bit and that ghost.
  An implementation that leaves the announcement to the adversary instantiates the relation as the full one, and the bit is then unconstrained (Definition~\ref{def:aba-ghost-free}).
  The content is the implementation's own, the two instantiations holding different ghosts and writing them at different transitions, which is why $G$, $\mathrm{ghostStep}$ and $\mathrm{ghostOutput}$ are parameters here as $M$, $S$ and the implementation's transitions are.
  Full construction: \leandecl{PLTS.ABA.Implementation.system}.
\end{definition}
